\documentclass[11pt,a4paper]{article}

% ---------- Packages ----------
\usepackage[margin=1in]{geometry}
\usepackage{graphicx}
\usepackage{float}        % for [H] exact figure/table placement
\usepackage{booktabs}     % nice tables
\usepackage{hyperref}
\usepackage{xcolor}
\usepackage{listings}     % code listings

\hypersetup{
    colorlinks=true,
    linkcolor=black,
    urlcolor=blue,
    citecolor=blue
}

% ---------- Code style ----------
\lstdefinestyle{mystyle}{
  basicstyle=\ttfamily\small,
  keywordstyle=\color{blue},
  commentstyle=\color{gray!70},
  stringstyle=\color{orange!70!black},
  numbers=left,
  numberstyle=\tiny\color{gray},
  stepnumber=1,
  numbersep=8pt,
  frame=single,
  breaklines=true,
  tabsize=2,
  showstringspaces=false
}
\lstset{style=mystyle, language=Python}

% ---------- Title ----------
\title{\textbf{Credit Card Fraud Detection}\\
\large Model Comparison using Logistic Regression, Decision Tree, and MLP (ANN)}
\author{Saurav Banerjee (Roll No: 16014123042) \\
Urvi Bhat (Roll No: 16014123059) \\
K. J. Somaiya College of Engineering}
\date{\today}

\begin{document}
\maketitle

\begin{abstract}
This mini-project builds and compares three models---Logistic Regression, Decision Tree, and a Multi-Layer Perceptron (ANN)---for credit card fraud detection. Using a Kaggle dataset, we implement preprocessing, train--test split, SMOTE for handling imbalance, training, and evaluation using Precision, Recall, F1-score, and ROC-AUC. Results show that the ANN attains the best overall performance.
\end{abstract}

\section{Problem Statement}
Credit card fraud is a major challenge for financial institutions. The goal is to classify transactions as \emph{fraudulent} or \emph{legitimate} using machine learning. We evaluate three models aligned with the syllabus:
\begin{itemize}
    \item Logistic Regression (baseline linear classifier),
    \item Decision Tree (interpretable, non-linear),
    \item Multi-Layer Perceptron (ANN).
\end{itemize}

\section{Dataset}
We use the ``Credit Card Fraud Detection Dataset 2023'' from Kaggle~\cite{dataset}.  
Features: anonymized variables $V1$--$V28$, transaction \texttt{Amount}, and target label \texttt{Class} (fraud=1, non-fraud=0).  
The \texttt{id} column was dropped since it is not predictive.

\section{Methodology}
The pipeline (Figure~\ref{fig:block}) includes:
\begin{enumerate}
    \item Preprocessing: scaling numerical features.
    \item Train/test split (80/20 stratified).
    \item SMOTE oversampling on training data to handle class imbalance.
    \item Model training: Logistic Regression, Decision Tree, ANN.
    \item Evaluation: Precision, Recall, F1-score, ROC-AUC.
\end{enumerate}

\begin{figure}[H]
    \centering
    \includegraphics[width=\linewidth]{Credit Card Fraud Detection.drawio.png}
    \caption{Workflow for Credit Card Fraud Detection.}
    \label{fig:block}
\end{figure}

\section{Code Snippet}
\begin{lstlisting}
logreg = ImbPipeline([
  ('pre', preprocess),
  ('smote', SMOTE(random_state=42)),
  ('clf', LogisticRegression(class_weight='balanced',
                             solver='liblinear',
                             max_iter=1000))
])

tree = ImbPipeline([
  ('pre', preprocess),
  ('smote', SMOTE(random_state=42)),
  ('clf', DecisionTreeClassifier(class_weight='balanced',
                                 random_state=42))
])

mlp = ImbPipeline([
  ('pre', preprocess),
  ('smote', SMOTE(random_state=42)),
  ('clf', MLPClassifier(hidden_layer_sizes=(64,32),
                        activation='relu',
                        early_stopping=True,
                        max_iter=120))
])
\end{lstlisting}

\section{Results}
Table~\ref{tab:metrics} shows the evaluation metrics of all three models on the test set.  
The ANN clearly achieves the highest values in Precision, Recall, F1-score, and ROC-AUC, 
indicating its superior performance. Logistic Regression performs slightly lower, while 
Decision Tree is very strong and close to ANN.

\begin{table}[H]
\centering
\caption{Model comparison on test set.}
\label{tab:metrics}
\begin{tabular}{lcccc}
\toprule
\textbf{Model} & \textbf{Precision} & \textbf{Recall} & \textbf{F1} & \textbf{ROC-AUC} \\
\midrule
Logistic Regression & 0.9773 & 0.9522 & 0.9646 & 0.9935 \\
Decision Tree       & 0.9968 & 0.9987 & 0.9978 & 0.9978 \\
MLP (ANN)           & \textbf{0.9991} & \textbf{1.0000} & \textbf{0.9996} & \textbf{0.99996} \\
\bottomrule
\end{tabular}
\end{table}

\noindent
From Table~\ref{tab:metrics}, we see that the ANN achieves almost perfect scores 
with Recall = 1.0, meaning no fraudulent transaction was missed. Decision Tree also 
performs very well but is slightly less accurate than ANN. Logistic Regression, being 
linear, is less powerful but still a good baseline.

\begin{figure}[H]
    \centering
    \includegraphics[width=0.8\linewidth]{roc_curves.png}
    \caption{ROC curves for the three models.}
    \label{fig:roc}
\end{figure}

\noindent
Figure~\ref{fig:roc} shows the ROC curves. The curve for ANN is almost touching 
the top-left corner, which indicates near-perfect classification ability. Decision Tree 
is also close, while Logistic Regression is slightly below but still strong.  
A higher ROC curve means better ability to distinguish fraud vs non-fraud.

\begin{figure}[H]
    \centering
    \includegraphics[width=0.6\linewidth]{confusion_matrix.png}
    \caption{Confusion matrix for best model (ANN).}
    \label{fig:cm}
\end{figure}

\noindent
Figure~\ref{fig:cm} presents the confusion matrix for ANN. It shows that almost 
all legitimate and fraudulent transactions were classified correctly. Only a very 
small number of false positives occurred (legitimate marked as fraud), while 
false negatives (fraud missed) are zero. This is desirable in fraud detection 
since missing frauds is more dangerous than raising a few false alarms.

\section{Future Scope}
This project can be extended in several directions:
\begin{itemize}
    \item Deploy the trained model as a REST API for real-time fraud detection in banking apps.
    \item Experiment with advanced models like Random Forests, Gradient Boosted Trees, or Transformer-based models.
    \item Perform cost-sensitive learning, where missing a fraud is penalized more than a false alarm.
    \item Explore feature importance to gain insights into which transaction attributes contribute most to fraud detection.
\end{itemize}

\section{Conclusion}
The MLP (ANN) outperformed Logistic Regression and Decision Tree, achieving nearly perfect recall and ROC-AUC. Decision Tree offers interpretability, while Logistic Regression provides a strong linear baseline. The project demonstrates that deep learning yields the highest fraud detection performance for this dataset.
\newpage
\begin{thebibliography}{9}
\bibitem{dataset}
Gayan Nuwani Nelgiriyewithana, ``Credit Card Fraud Detection Dataset 2023,'' Kaggle, 2023. Available: \url{https://www.kaggle.com/datasets/nelgiriyewithana/credit-card-fraud-detection-dataset-2023}

\bibitem{sklearn}
Pedregosa, F. et al., ``Scikit-learn: Machine Learning in Python,'' JMLR, vol. 12, pp. 2825–2830, 2011.

\bibitem{imblearn}
Lemaître, G., Nogueira, F., and Aridas, C.K., ``Imbalanced-learn: A Python Toolbox to Tackle the Curse of Imbalanced Datasets,'' JMLR, vol. 18, no. 17, pp. 1–5, 2017.
\end{thebibliography}

\end{document}
